\documentstyle[%


righttag%


,11pt%


,amssymb%


,russian%               remove in Eng.


,izvru%                 Set izven% in Eng.


% ,izvru-ks             for brief communications


]{amsart}





%       Math definitions





\newcommand{\dstyle}{\displaystyle}


\newcommand{\down}[2]{\underset{#1}{#2}{}}


\newcommand{\up}[2]{\overset{#1}{#2}{}}


%\newcommand{\wt}{\widetilde}


%\newcommand{\ov}{\overline}


%\newcommand{\wh}{\widehat}


\newcommand{\tts}[1]{}





%\hoffset=-15mm%                  For Eng.





\setcounter{page}{29}


\god{1998}


\nomer{8~(435)} %                        Remove in Eng.


%\nomer{Vol.\,42, No.\,8}%               Set in Eng.


%\str{\pageref{begin}--\pageref{end}}%  Set in Eng.


\udk{514.762}


\author{\it Р.Х.\,Ибрагимова}


% NB:   ^^ remove in Eng.


\title{Движения на касательных расслоениях,\\


сохраняющие ортогональную \\ и касательную структуры}





\date{}


% \pagestyle{myheadings}%         Set in Eng.





\pagestyle{plain} %              Remove in Eng.


\begin{document}





% \thispagestyle{plain}%          Set in Eng.





\maketitle


\label{begin}





В \cite{1} были изучены автоморфизмы произвольных расслоенных


пространств, относительно которых инвариантна $\pi$-структура; найдено


строение тензоров кривизны и кручения для расслоенного многообразия,


допускающего максимальную группу автоморфизмов; рассмотрены автоморфизмы


более специального типа, при которых сохраняется также и заданный объект


линейной связности.


В \cite{2} рассматриваются три группы движений в касательном расслоении


со специальной метрикой типа Сасаки: движения, сохраняющие расслоенную


структуру; движения, состоящие из продолженных преобразований, и


произвольные движения. Найдены максимальные размерности этих


групп движений.





В данной работе изучаются движения на касательном расслоении с


произвольной римановой метрикой, не вырождающейся на слоях, сохраняющие


ортогональную и касательную структуры, в предположении, что метрическая


связность является канонической, а распределения $H$ и $V$\;


$J$-изометричны. Вычисления проводятся относительно неголономного поля


реперов. При этом используется аппарат теории автоморфизмов в


расслоенных пространствах, построенный в \cite{1}.





Рассмотрим на касательном расслоении $TM$ риманову метрику произвольной


сигнатуры, не вырожденную на слоях. Задав любое (вообще говоря,


неголономное) адаптированное к $\pi$-структуре


поле реперов $\{e_A\}$ и сопряженное ему поле кореперов


$\{\theta^A\}$: $\theta^A(e_B)=\delta^A_B$, будем иметь


\begin{equation}


d\sigma^2=G_{AB}\theta^A\theta^B,


\end{equation}


где


$$


G_{AB}=


\begin{pmatrix}


G_1 & G_4 \\ G_3 & G_2


\end{pmatrix}


,\quad G_1=(G_{ij}),\quad G_2=(G_{\ov i\,\ov j}),\quad


G_3=(G_{\ov i j}),\quad G_4=(G_{i\ov j}).


$$


Пусть на $TM$ задано произвольное векторное поле $v=v^A e_A$,


определяющее инфинитезимальное преобразование


\begin{equation}


X^{A'}=X^A+v^A\tau.


\end{equation}


Преобразование (2) определяет движение, если оно сохраняет метрику,


т.\,е. векторное поле $v$ удовлетворяет уравнениям Киллинга


\begin{equation}


\down{v}{D}G_{AB}=0.


\end{equation}


Разобьем их на три группы


\begin{equation}


\down{v}{D}G_{ij}=0,\quad \down{v}{D}G_{\ov i\,\ov j}=0,\quad


\down{v}{D}G_{\ov i j}=0.


\end{equation}





Пусть $\nabla$ --- метрическая связность с заданным тензором кручения


\cite{3}, т.\,е. $\nabla$ удовлетворяет условию $\nabla_C G_{AB}=0$. Тогда


уравнения Киллинга (3) принимают вид


$$


2v_{(AB)}=v_{AB}+v_{BA}=0,


$$


где $v_{BA}=\nabla_A v_B+S_{BAC}v^C$, $S$ --- тензор кручения связности,


$v_B=v^A G_{AB}$. Таким образом,


\begin{equation}


\nabla_B v_A=v_{AB}+S_{ACB}v^C.


\end{equation}


Тензор кривизны касательного расслоения $(TM,G)$


с произвольной римановой


метрикой $G$, допускающего максимальную группу движений, обладает


структурой


\begin{equation}


K^A_{BCD}=\theta^A_{BCD}+T^A_{BCD},


\end{equation}


где


$\theta^A_{BCD}=\frac{Q}{N(N-1)}(\delta^A_D G_{BC}-


\delta^A_C G_{BD})$


--- тензор кривизны римановой


связности,


\begin{equation}


\begin{split}


& T^A_{BCD}=\nabla_C P^A_{BD}-\nabla_D P^A_{BC}+


P^A_{MC}P^M_{BD}-P^A_{MD}P^M_{BC}-\tfrac{1}{2}P^A_{BM}


(P^M_{CD}-P^M_{DC}), \\


& P_{ABC}=S_{ABC}+S_{BCA}-S_{CAB}, \\


& Q_{BC}=Q^A_{BCA}=K_{BC}+\nabla_F P^F_{BC}-


\nabla_C P_B-P_F P^F_{BC}+P^F_{BM}P^M_{CF},\quad


P_B=P^M_{BM},


\end{split}


\end{equation}


$K_{BC}$ --- тензор Риччи.





Рассмотрим на касательном расслоении $TM$ с произвольной метрикой (1)


движения, относительно которых инвариантна заданная ортогональная


$\pi$-структура \cite{3}.





Выберем в качестве горизонтального распределения распределение ${}^\bot H$,


ортогональное слоям, т.\,е.


${}^\bot H\bot V$. В этом случае на $TM$ определяется $\pi$-структура,


которую будем называть \cite{3}


ортогональной. Задание этой структуры сводится к заданию аффинора $F$,


матрица которого в произвольном адаптированном к $\pi$-структуре


поле реперов $\{e_A\}$ имеет вид


$$


F=


\begin{pmatrix}


E_1 & 0 \\ 0 & -E_2


\end{pmatrix}


,\quad E_1=(\delta^i_j),\quad E_2=(\delta^{\ov i}_{\ov j})=


(\delta^i_j).


$$


Аффинор $F$ сохраняет скалярное произведение, т.\,е. $(Fu,Fv)=(u,v)$,


где $u$, $v$ --- произвольные векторные поля на $TM$, и метрический тензор


в адаптированном репере $\{e_A\}$ имеет компоненты


\begin{equation}


(G_{AB})=


\begin{pmatrix}


G_1 & 0 \\ 0 & G_2


\end{pmatrix}


,


\end{equation}


т.\,е. $G_{i\ov j}=0$. Как известно \cite{1}, преобразования (2)


сохраняют слои тогда и только тогда, когда векторное поле $v$ является


проектируемым, т.\,е. в адаптированном репере $v^i=v^i(x^k)$.





\begin{thm}


Если проектируемое векторное поле $v$ является движением метрики $(8)$,


то оно сохраняет горизонтальное распределение ${}^\bot H$.


\end{thm}





\begin{pf}


Пусть проектируемое в.\,п. $v$ есть движение, ${}^\bot H$ ---


горизонтальное


распределение, ортогональное слоям. Тогда в адаптированном репере


$G_{i\ov j}=0$ и


второе из уравнений (4) дает


$G_{\ov r\ov j}v^{\ov r}_{\ov i}=0$, откуда в силу невырожденности блока


$G_2=(G_{\ov r\ov j})$ следует $v^{\ov r}_i=0$,


что равносильно \cite{4} условию


$\down{v}{D}F^A_B=0$.


\end{pf}





Для инвариантности заданной $\pi$-структуры относительно преобразования (2)


необходимо и достаточно выполнения условия \cite{1}


\begin{equation}


\down{v}{D}F^A_B=0.


\end{equation}


Это условие в адаптированном репере для проектируемого векторного поля


$v$ запишется в виде


$$


\nabla_v F^A_B-F^C_B v^A_C+F^A_C v^C_B=0.


$$


Если связность вполне приводима, т.\,е.


$\nabla_v F^A_B=0$, то $F^A_C v^C_B=F^C_B v^A_C$. Тогда в


адаптированном репере матрица $v_{AB}$ имеет вид


\begin{equation}


v_{AB}=


\begin{pmatrix}


v_{ij} & 0 \\ 0 & v_{\ov i\,\ov j}


\end{pmatrix}


,


\end{equation}


где блоки кососимметричны, что эквивалентно соотношениям


$v_{(ij)}=0$, $v_{(\ov i\,\ov j)}=0$.


Поэтому наряду с $N=2n$ неизвестными компонентами


векторного поля $v^A(v^i,v^{\ov i})$ рассмотрим также


$n(n-1)/2+n(n-1)/2=n(n-1)$ кососимметричных компонент


$v_{ij}$, $v_{\ov i\,\ov j}$. Порядок группы движений $(TM,G)$ с метрикой


(8), относительно которых инвариантна


заданная $\pi$-структура, не превосходит, следовательно, числа


$r=n(n-1)+2n=n(n+1)$.





Итак, справедлива





\begin{thm}


Порядок группы проектируемых движений касательного расслоения $(TM,G)$


с метрикой $(8)$ не превосходит числа $r=n(n+1)$.


\end{thm}





Уравнения (9) являются следствием уравнений Киллинга для метрики (8).


Уравнения Киллинга (4) в силу соотношения (9) принимают вид


\begin{equation}


\down{v}{D}G_{ij}=0,\quad \down{v}{D}G_{\ov i\,\ov j}=0.


\end{equation}


Что касается уравнений $\down{v}{D}G_{i\ov j}=0$,


то они удовлетворяются тождественно. Следовательно, имеет место





\begin{thm}


Для того чтобы однопараметрическая группа автоморфизмов ортогональной


$\pi$-структуры была движением, необходимо и достаточно, чтобы


соответствующее векторное поле $v$ удовлетворяло дифференциальным уравнениям


$(11)$.


\end{thm}





В дальнейших рассуждениях для упрощения вычислений ограничимся


рассмотрением канонической связности, удовлетворяющей условиям \cite{3}


$$


\nabla F=0,\quad \nabla J=0,\quad P_{ijk}=0,\quad


P_{\ov i\,\ov j\,\ov k}=0.


$$


Выясним строение тензоров кривизны и кручения этой связности. Так как


связность каноническая, то для формы кривизны $\Omega$ этой связности


имеем $\Omega^i_{\ov j}=0$, $\Omega^{\ov i}_j=0$,


$\Omega^i_j=\Omega^{\ov i}_{\ov j}$, т.\,е. тензор


кривизны обладает отличными от нуля компонентами


\begin{equation}


K^m_{jei}=K^{\ov m}_{\ov j ei},\quad


K^m_{je\ov i}=K^{\ov m}_{\ov j e\ov i},\quad


K^m_{j\ov e\ov i}=K^{\ov m}_{\ov j\ov e\ov i}.


\end{equation}


Учитывая эти соотношения, согласно (6) для тензора кривизны $K^A_{BCD}$


канонической связности получаем отличные от нуля компоненты


\begin{equation}


K^m_{jei}=Q^m_{jei}+T^m_{jei},\quad


K^m_{je\ov i}=Q^m_{je\ov i}+T^m_{je\ov i},\quad


K^m_{j\ov e\ov i}=Q^m_{j\ov e\ov i}+T^m_{j\ov e\ov i},


\end{equation}


где согласно (7)


\begin{align*}


T^m_{jei}&=P^m_{\ov r j}P^{\ov r}_{ei}-


P^m_{\ov r i}P^{\ov r}_{j\ov e}-P^m_{j\ov r}


(P^{\ov r}_{ei}-P^{\ov r}_{ie}), \\


T^m_{je\ov i}&=\nabla_e P^m_{j\ov i}-P^m_{re}P^r_{j\ov i}-


P^m_{\ov r e}P^{\ov r}_{j\ov i}-


P^m_{\ov r i}P^{\ov r}_{je}-\tfrac{1}{2}P^m_{j\ov r}


(P^{\ov r}_{e\ov i}-P^{\ov r}_{\ov i e}), \\


T^m_{j\ov e\ov i}&=\nabla_{\ov e}P^m_{j\ov i}-


\nabla_{\ov i}P^m_{j\ov e}+P^m_{c\ov e}P^c_{j\ov i}-


P^m_{c\ov i}P^c_{j\ov e}, \\


Q^m_{jei}&=\tfrac{Q}{n(n-1)}(\delta^m_i G_{je}-


\delta^m_e G_{ji}).


\end{align*}





Таким образом, доказана





\begin{thm}


Тензор кривизны канонической связности касательного расслоения $(TM,G)$


с метрикой $(8)$, допускающего максимальную группу движений, сохраняющих


ортогональную $\pi$-структуру, имеет вид $(13)$.


\end{thm}





Тензор кручения канонической связности кроме заданных условий


$P_{ijk}=0$, $P_{\ov i\,\ov j\,\ov k}=0$


удовлетворяет условиям, вытекающим из соотношений (12),


\begin{equation}


\begin{split}


& T^m_{\ov j\ov e\ov i}=0,\quad T^{\ov m}_{jei}=0,\quad


T^{\ov m}_{\ov j ei}=0,\quad T^{\ov m}_{j\ov e\ov i}=0,\quad


T^m_{\ov j e\ov i}=Q^m_{\ov j e\ov i}, \\


& T^{\ov m}_{j e\ov i}=-Q^{\ov m}_{j e\ov i},\quad


T^{\ov m}_{\ov j e\ov i}-T^m_{jei}=Q^m_{jei},\quad


\nabla_r T^m_{\ov j\ov e i}=0, \\


& T^{\ov m}_{\ov j e\ov i}=T^{\ov m}_{j e\ov i},\quad


T^m_{j\ov e\ov i}-T^{\ov m}_{\ov j\ov e\ov i}=


Q^{\ov m}_{\ov j\ov e\ov i},\quad


\nabla_{\ov r}T^m_{\ov j\ov e i}=0, \\


& P^{\ov k}_{\ov r m}(P^{\ov r}_{ke}\delta^m_i-


P^{\ov r}_{ki}\delta^m_e)=0,\quad


P^r_{\ov k\ov m}(P^k_{r\ov i}\delta^{\ov m}_{\ov e}-


P^k_{r\ov e}\delta^{\ov m}_{\ov i}=0,


\end{split}


\end{equation}


где различные компоненты тензора $T^A_{BCD}$


вычисляются согласно (7), а


$$


Q^{\ov m}_{\ov j\ov e\ov i}=\frac{1}{n-1}


(\delta^{\ov m}_{\ov i}\wt Q_{\ov j\ov e}-\delta^{\ov m}_{\ov e}


\wt Q_{\ov j\,\ov i}),\quad


\wt Q_{\ov j\ov e}=Q^{\ov m}_{\ov j\ov e\,\ov m}.


$$





Итак, справедлива





\begin{thm}


Тензор кручения канонической связности касательного расслоения


$(TM,G)$ с


метрикой $(8)$, допускающего максимальную группу движений, сохраняющих


ортогональную $\pi$-структуру, удовлетворяет условиям $(14)$.


\end{thm}





В частности, если каноническая связность $\nabla$ риманова, то $P=0$,


$T=0$. Тогда $K_{AB}=Q_{AB}=0$, т.\,е. $K^D_{ABC}=0$. В этом случае


\cite{3} касательное расслоение $TM$,


допускающее вполне приводимую риманову связность, тривиально и поэтому


является прямым произведением двух римановых пространств с метриками


$G_{ij}(x^k)$, $G_{\ov i\,\ov j}(x^k)$. Таким образом, имеет место


следующая





\begin{thm}


Тензор кривизны канонической римановой связности касательного расслоения


$(TM,G)$ с метрикой $(8)$, допускающего максимальную группу движений,


сохраняющих ортогональную $\pi$-структуру, равен нулю.


\end{thm}





Особенностью касательного расслоения является наличие изоморфизма


векторных подрасслоений


$$


J:{}^\bot H\to V,


$$


определяемое аффинором касательной структуры $J$ ($J^2=0$), матрицей


которого в индуцированных координатах будет


$$


J=


\begin{pmatrix}


0 & 0 \\ E & 0


\end{pmatrix}


,\quad E=(\delta^i_j);\quad


\text{Ker\,}(J)=\text{Im\,}(J)=V(TM).


$$





Рассмотрим на касательном расслоении $(TM,G)$ с метрикой (8) движения,


сохраняющие кроме того и касательную структуру. Необходимым и


достаточным условием инвариантности касательной структуры относительно


преобразования (2) является равенство


\begin{equation}


\down{v}{D}J^A_B=0.


\end{equation}


Для проектируемого векторного поля $v$ условия (15) принимают вид


$$


\nabla_v J^A_B+J^A_C v^C_B-J^C_B v^A_C=0.


$$


Так как связность каноническая, то $\nabla_v J^A_B=0$. Поэтому


$J^A_C v^C_B=J^C_B v^A_C$. Тогда в


адаптированном репере $v_{ij}=v_{\ov i\,\ov j}$


и матрица (10) принимает вид


$$


v_{AB}=


\begin{pmatrix}


v_{ij} & 0 \\ 0 & v_{ij}


\end{pmatrix}


,


$$


где блоки кососимметричны в силу уравнений $v_{(AB)}=0$, которые


эквивалентны


соотношению $v_{(ij)}=0$ и поэтому $v_{ij}=v_{[ij]}$. Тогда наряду с


$N=2n$ неизвестными компонентами векторного поля $v^A(v^i,v^{\ov i})$


рассмотрим еще $n(n-1)/2$ кососимметричных компонент $v_{ij}$. Таким


образом, порядок группы движений на касательном расслоении $(TM,G)$


с метрикой (8),


относительно которых инвариантны ортогональная и касательная структуры,


не превосходит числа


\begin{equation}


r=2n+n(n-1)/2=(n^2+3n)/2.


\end{equation}





Итак, имеет место





\begin{thm}


Порядок группы движений на касательном расслоении $(TM,G)$ с метрикой


$(8)$, сохраняющих $\pi$-структуру и касательную структуру $J$, не


превосходит числа $r=(n^2+3n)/2$.


\end{thm}





Рассмотрим пример. Если касательное расслоение $TM$ является евклидовым


пространством $E^{2n}$, то для евклидовой метрики $G_{AB}=\delta_{AB}$


уравнения Киллинга


(3) имеют вид $\partial_A v_B+\partial_B v_A=0$. В этом случае задача


нахождения движений на $(TM,G)$, сохраняющих


ортогональную $\pi$-структуру и касательную структуру $J$,


сводится к интегрированию системы уравнений


$$


\partial_{\ov i}v_{\ov j}=\partial_i v_j,\quad


\partial_i v_j+\partial_j v_i=0,\quad


\partial_{\ov i}v_j=0,\quad \partial_i v_{\ov j}=0.


$$


Интегрируя, получим


$$


v_{\ov i}=c_{ij}x^j+c_{\ov i},\quad c_{ij}=-c_{ji},\quad


v_i=c_{ij}x^j+c_i,\quad v_i=v_i(x^k),\quad


v_{\ov i}=v_{\ov i}(x^{\ov k}).


$$


Таким образом, решение данной системы зависит от $2n$ произвольных


постоянных $c_i$, $c_{\ov i}$ и\linebreak $n(n-1)/2$ независимых


элементов кососимметричной матрицы


$c_{ij}$, т.\,е. максимальный порядок группы таких движений точно


равен числу в (16).





Выясним строение тензоров кривизны и кручения канонической связности.


Учитывая (12), из соотношений (6) получим для тензора кривизны


$K^A_{BCD}$ канонической связности отличные от нуля


компоненты


\begin{equation}


K^m_{jei}=Q^m_{jei}+T^m_{jei},\quad


K^m_{j\ov e i}=Q^m_{j\ov e i}+T^m_{j\ov e i},\quad


K^m_{j\ov e\ov i}=Q^m_{j\ov e\ov i}+T^m_{j\ov e\ov i},


\end{equation}


где $T^m_{jei}$, $T^m_{j\ov e i}$, $T^m_{j\ov e\ov i}$


выражаются согласно (7),


\begin{align*}


Q^m_{jei}&=\frac{Q}{n(n-1)}(\delta^m_i G_{je}-


\delta^m_e G_{ji}), \\


%


Q^m_{j\ov e i}&=\frac{1}{n}\{(\delta^m_i\wt{\wt Q}_{je}-


\delta^m_e\wh Q_{ji})+(T^{\ov m}_{\ov j\ov e i}-


T_{\ov j\ov e i}{}^{\ov m}_{\dstyle\cdot}-


T_{\ov j i}{}^{\ov m}_{\dstyle\cdot}


{}_{\ov e})-(T_{j\ov e i}{}^m_{\dstyle\cdot}-T^m_{j i\ov e}-


T_{jei}{}^{\ov m}_{\dstyle\cdot})\}, \\


%


\wt{\wt Q}_{AB}&=Q^m_{AB m}, \\


%


Q^m_{j\ov e\ov i}&=\frac{1}{n}\{(\delta^{\ov m}_{\ov i}


\wh Q_{j\ov e}-\delta^{\ov m}_{\ov e}\wh Q_{j\ov i})+


(T_{\ov j\ov e}{}^{\ov m}_{\dstyle\cdot}{}_i-T^{\ov m}_{\ov j\,\ov i\ov e}+


T_{\ov j\,\ov i\ov e}{}^m_{\dstyle\cdot}) -


(T_{ij\ov e}{}^{\ov m}_{\dstyle\cdot}-


T^m_{j\ov i\ov e}+T_{je}{}^{\ov m}_{\dstyle\cdot}{}_{\ov i})\}, \\


%


\wh Q_{AB}&=Q^m_{AB\ov m}.


\end{align*}





Итак, справедлива





\begin{thm}


Тензор кривизны канонической связности касательного расслоения $(TM,G)$


с метрикой $(8)$, допускающего максимальную группу движений, сохраняющих


ортогональную $\pi$-структуру и касательную структуру, имеет вид $(17)$.


\end{thm}





Тензор кручения канонической


связности кроме заданных условий $P_{ijk}=0$,


$P_{\ov i\,\ov j\,\ov k}=0$ удовлетворяет условиям,


вытекающим из соотношений (12),


\begin{gather}


T^{\ov m}_{jei}=0, \quad T^{\ov m}_{je\ov i}=0,\quad


T^{\ov m}_{j\ov e\ov i}=0,\quad T^m_{\ov j ei}=0,\quad


T^m_{\ov j e\ov i}=0, \notag\\


%


T^m_{\ov j\ov e\ov i}=0,\quad \wh T_{\ov i\ov e}=


\wt T_{i\ov e},\quad \wh T_{\ov e k}=\up{*}{T}_{k\ov e},\quad


\wh T_{ki}=\wt{\wt T}_{\ov k\,\ov i},\quad


\wh T_{\ov e i}=\wh T_{\ov i e}, \\


%


\up{*}{T}_{ei}=\wt{\wt T}_{\ov e i},\quad


\wt{\wt T}_{\ov i e}=\wt{\wt T}_{\ov e i},\quad


\wt T_{e\ov i}=\wt T_{i\ov e},\quad


\wh T_{\ov e\ov i}=\wh T_{\ov i\ov e},\notag


\end{gather}


где $\wt T_{AB}=T^{\ov m}_{AB\ov m}$,


$\wt{\wt T}_{AB}=T^m_{ABm}$,


$\wh T_{AB}=T^m_{AB\ov m}$,


$\up{*}{T}_{AB}=T^{\ov m}_{ABm}$, а компоненты


$T^A_{BCD}$


выражаются согласно (7). Отсюда следует





\begin{thm}


Тензор кручения канонической связности на касательном расслоении


$(TM,G)$ с метрикой $(8)$, допускающем максимальную группу движений,


сохраняющих ортогональную $\pi$-структуру и касательную структуру,


удовлетворяет условиям $(18)$.


\end{thm}





\begin{thebibliography}{9}





\bibitem{1}


Шапуков Б.Н. {\em Автоморфизмы расслоенных пространств} // Тр.


Геометрич. семин. -- Казань, 1982, вып.\,14. -- С.\,97--108.


\bibitem{2}


Паньженский В.И. {\em О движениях в касательном расслоении с метрикой


Сасаки}. -- Пенз. гос. пед. ин-т. -- Пенза, 1989. -- 10~с. -- Деп. в


ВИНИТИ 10.02.89, \No\,1194-B89.


\bibitem{3}


Ибрагимова Р.Х., Шапуков Б.Н. {\em Метрические расслоения и некоторые их


приложения} // Тр. Геометрич. семин. -- Казань, 1990, вып.\,20. --


С.\,44--58.


\bibitem{4}


Ибрагимова Р.Х. {\em Движения на касательных расслоениях со специальной


метрикой} // Дифференц. геометрия. -- Саратов, 1985, вып.\,8. --


С.\,17--22.





\end{thebibliography}





\vskip 2cm





\post{\it Алмаатинский университет}{\it Поступила}


\post{\it $($Казахстан$)$}{22.12.1995}





\label{end}


\end{document}


